%%
%% Joule: Energy-Efficient Agent Task Execution
%% ACM sigconf — MLSys '26
%%
\documentclass[sigconf,authordraft]{acmart}

\AtBeginDocument{%
  \providecommand\BibTeX{{%
    Bib\TeX}}}

\setcopyright{acmlicensed}
\copyrightyear{2026}
\acmYear{2026}
\acmDOI{XXXXXXX.XXXXXXX}

\acmConference[MLSys '26]{Machine Learning and Systems}{2026}{USA}
\acmISBN{978-1-4503-XXXX-X/2026/06}

\usepackage{enumitem}
\usepackage{booktabs}
\usepackage{multirow}
\usepackage{microtype}

\begin{document}

\title{Joule: Energy-Aware Orchestration for Multi-Step LLM Agent Workflows\\
via Adaptive Context Selection, Routing, and Budget Enforcement}

\author{Aagam Bothara}
\email{aboth@uic.edu}
\affiliation{%
  \institution{University of Illinois Chicago}
  \city{Chicago}
  \state{Illinois}
  \country{USA}
}

\renewcommand{\shortauthors}{Bothara}

% ─────────────────────────────────────────────────────────────────────
\begin{abstract}
Multi-step LLM agent workflows suffer from two sources of avoidable overhead.
First, \emph{context propagation waste}: frameworks naively replay all accumulated step results into every downstream context window, even when most are irrelevant to the current step.
Second, \emph{planning overhead}: a na\"{i}ve pipeline issues four sequential LLM calls (specification, classification, planning, critique) before any tool use begins, multiplying token consumption for every task.
Neither problem is addressed by existing frameworks; neither CrewAI nor LangChain models the dependency structure of execution plans or minimizes planning call overhead.

We present \textbf{Joule}, an execution-aware runtime for multi-step LLM agents that treats context management and call efficiency as first-class runtime concerns.
Joule's primary contribution is \emph{confidence-aware dependency context selection}: it models each execution plan as a DAG, computes the transitive closure of relevant step results at synthesis time, and scores each candidate-for-pruning step using three runtime signals (data-flow taint tracking, lexical relevance, output richness) to decide whether pruning is safe.
An \emph{adaptive controller} raises the pruning conservatism threshold dynamically based on plan length, token budget pressure, and observed planner dependency under-declaration rate, preventing the quality regressions that afflict binary pruning.

Supporting mechanisms include: (1)~\emph{SLM-first routing} with a five-rule cascade that defaults to small models; (2)~\emph{adaptive prompt optimization} that unifies four planning calls into one---reducing prompt tokens by 80\% and LLM calls by 71\% vs.\ a na\"{i}ve pipeline; and (3)~\emph{7-dimensional budget enforcement} that provides a hard per-task resource ceiling available in no prior framework.

We evaluate Joule against CrewAI and LangChain on a 50-task benchmark (3 runs per task) spanning five categories.
On \emph{research} and \emph{multi-step} tasks---where multi-step planning adds genuine value---Joule reduces token-weighted inference cost\footnote{We report \emph{token-weighted energy proxy} (Wh) throughout: $E = (n_{\text{in}} \cdot \alpha_{\text{in}} + n_{\text{out}} \cdot \alpha_{\text{out}}) / 10^6$, not direct hardware measurement.} by 28\% and 30\% vs.\ CrewAI, while achieving \textbf{3.9$\times$} and \textbf{2.8$\times$} lower latency.
The confidence-aware adaptive controller eliminates the quality regression seen with binary pruning: multi-step quality fully recovers to parity with no-pruning (4.40 vs.\ 4.40), and overall quality improves from 4.49 to 4.61.
Prompt optimization is the dominant energy mechanism: disabling it multiplies energy proxy by \textbf{4.0$\times$} and LLM calls from 1.5 to 5.3 per task.
Budget enforcement provides zero resource violations across 150 executions.
\end{abstract}

\begin{CCSXML}
<ccs2012>
  <concept>
    <concept_id>10010520.10010553.10010562</concept_id>
    <concept_desc>Computer systems organization~Autonomous agents</concept_desc>
    <concept_significance>500</concept_significance>
  </concept>
  <concept>
    <concept_id>10010520.10010575.10010577</concept_id>
    <concept_desc>Computer systems organization~Energy aware systems</concept_desc>
    <concept_significance>300</concept_significance>
  </concept>
  <concept>
    <concept_id>10010147.10010178.10010187</concept_id>
    <concept_desc>Computing methodologies~Intelligent agents</concept_desc>
    <concept_significance>100</concept_significance>
  </concept>
</ccs2012>
\end{CCSXML}

\ccsdesc[500]{Computer systems organization~Autonomous agents}
\ccsdesc[300]{Computer systems organization~Energy aware systems}
\ccsdesc[100]{Computing methodologies~Intelligent agents}

\keywords{LLM agents, energy efficiency, context selection, dependency-aware pruning, adaptive routing, budget enforcement, agent runtime}

\maketitle

% ─────────────────────────────────────────────────────────────────────
\section{Introduction}
\label{sec:intro}

Agent frameworks that decompose tasks into sequences of LLM calls, tool invocations, and structured reasoning steps have become the standard pattern for automating complex workflows.
Yet the dominant frameworks---CrewAI, LangChain, AutoGen---were designed for \emph{correctness and composability}, not \emph{resource efficiency}.
Two structural inefficiencies compound as workflows grow:

\paragraph{Context propagation waste.}
When an agent executes a 5-step plan, the standard pattern is to pass \emph{all} prior step results as context to each synthesis or recovery call.
If step~5 only depends on step~3, the outputs of steps~1, 2, and~4 consume input tokens for no benefit.
No framework we surveyed models the dependency structure of execution plans or uses it to select relevant context.
The waste is invisible and compounds with plan length.

\paragraph{Planning call overhead.}
A na\"{i}ve planning pipeline issues four sequential LLM round-trips before tool use begins: (1)~extract a task specification, (2)~classify complexity, (3)~generate a step plan, (4)~critique the plan.
Each call carries a full system prompt with all tool descriptions ($\sim$2,900 tokens for a typical agent).
Multiplied across tasks, this overhead is the largest avoidable source of token consumption in the pipeline.

This paper presents \textbf{Joule}, an execution-aware runtime that addresses both inefficiencies.
Its central mechanism is \emph{confidence-aware dependency context selection}: rather than passing all step results or naively pruning by graph reachability alone, Joule scores each candidate-for-pruning step using runtime signals to determine whether pruning is safe, and defers to an adaptive controller that raises conservatism based on observed planning uncertainty.

\paragraph{Scope.}
Joule is designed for the workloads where agent frameworks provide value: multi-step research, code generation, and automation tasks requiring 2+ tool calls.
For pure question answering and text summarization---which need no planning machinery---direct LLM calls are more efficient.
We include these categories in our evaluation to be transparent, reporting them as expected-overhead cases with honest explanation.

\paragraph{Contributions.}
\begin{itemize}[leftmargin=2em,itemsep=2pt]
  \item A formal energy proxy model for multi-step agent workflows (\S\ref{sec:energy}), grounding token counts in energy-cost terms.
  \item \emph{Confidence-aware dependency context selection}: DAG reachability + three runtime signals (taint tracking, lexical relevance, output richness) + an adaptive threshold controller, eliminating the quality regressions of binary pruning (\S\ref{sec:pruning}).
  \item An \emph{adaptive controller} that adjusts the pruning conservatism threshold based on plan length, budget pressure, and dependency under-declaration rate (\S\ref{sec:adaptive}).
  \item \emph{Adaptive prompt optimization}: unified planning (4 calls $\to$ 1), slim prompts, direct-answer shortcut---reducing prompt tokens by 80\% and LLM calls by 71\% (\S\ref{sec:prompt-opt}).
  \item A five-rule SLM-first routing cascade with bidirectional budget coupling (\S\ref{sec:routing}).
  \item 7-dimensional budget enforcement providing hard per-task resource guarantees (\S\ref{sec:budget}).
  \item An empirical evaluation on 150 benchmark runs against CrewAI and LangChain with ablation studies and an honest per-category breakdown (\S\ref{sec:eval}).
\end{itemize}

% ─────────────────────────────────────────────────────────────────────
\section{Background}
\label{sec:background}

\subsection{Energy Proxy Model}
\label{sec:energy}

The energy consumed by a single LLM inference call is a function of token counts and per-model energy coefficients.
Luccioni et al.~\cite{luccioni2023power} show that large frontier models consume 4--10$\times$ more energy per token than small models, with output generation 3--5$\times$ more expensive than input prefill due to autoregressive decoding's memory-bandwidth bottleneck.

We adopt the following \emph{energy proxy} throughout:
\[
  E_{\text{call}} = \frac{n_{\text{in}} \cdot \alpha_{\text{in}} + n_{\text{out}} \cdot \alpha_{\text{out}}}{10^6}\ \text{Wh}
\]
where $n_{\text{in}}$, $n_{\text{out}}$ are input and output token counts and $\alpha_{\text{in}}$, $\alpha_{\text{out}}$ are per-model coefficients in Wh per million tokens (Table~\ref{tab:energy-tiers}).
This is a \emph{proxy}, not a direct hardware measurement; it models inference cost as proportional to token consumption weighted by model efficiency.
Per-task cost sums over all LLM calls issued during execution.
We use ``energy proxy'' and ``Wh'' interchangeably in results tables, with the understanding that actual hardware efficiency depends on deployment hardware, batching, and utilization.

\begin{table}[h]
\centering
\caption{Energy proxy coefficients (Wh/M tokens) for evaluation models.
Derived from Luccioni et al.~\cite{luccioni2023power} and vendor disclosures~\cite{openai2024efficiency,anthropic2024efficiency}.}
\label{tab:energy-tiers}
\begin{tabular}{llcc}
\toprule
Model & Tier & $\alpha_{\text{in}}$ & $\alpha_{\text{out}}$ \\
\midrule
\texttt{gpt-4o-mini}         & SLM & 0.30 & 1.20 \\
\texttt{gpt-4o}              & LLM & 1.50 & 5.00 \\
\texttt{claude-haiku-4.5}    & SLM & 0.40 & 1.50 \\
\texttt{claude-sonnet-4.6}   & LLM & 1.20 & 4.50 \\
\texttt{llama3.2:3b} (local) & SLM & 0.15 & 0.60 \\
\bottomrule
\end{tabular}
\end{table}

\subsection{The Two Structural Inefficiencies}

\paragraph{Context accumulation.}
In a 5-step plan, the synthesis call typically receives all 5 step results as context.
If the final answer only depends on the output of steps~3 and~5, steps~1, 2, and~4 contribute only noise---consuming input tokens for no informational benefit.
The energy proxy wasted is proportional to the number of irrelevant step outputs $\times$ their average token length $\times$ $\alpha_{\text{in}}$.
For long plans with sparse dependency graphs, this is the dominant source of waste.

\paragraph{Planning overhead.}
A na\"{i}ve 4-call planning pipeline incurs $\sim$4 $\times$ 2,900 = 11,600 prompt tokens before tool execution begins.
Our ablation (\S\ref{sec:ablation}) confirms this: the \emph{no-prompt-opt} configuration averages 11,023 prompt tokens and 5.3 LLM calls per task; full Joule averages 2,162 tokens and 1.5 calls.

\subsection{Existing Frameworks}
CrewAI~\cite{crewai}, LangChain~\cite{langchain}, and AutoGen~\cite{autogen} address task decomposition and composability but provide no context selection, no prompt call minimization, and no resource budget enforcement.
LangChain's ReAct~\cite{yao2022react} uses an iterative observe-think-act loop that typically issues multiple sequential LLM calls per task.
This motivates Joule's design.

% ─────────────────────────────────────────────────────────────────────
\section{System Design}
\label{sec:system}

\subsection{Confidence-Aware Dependency Context Selection}
\label{sec:pruning}

This is Joule's primary contribution.
The goal is to select, at synthesis time, exactly the step results that are relevant to the final answer---no more, no fewer.

\subsubsection{DAG-Based Plan Representation}

During planning, each step declares the 0-based indices of prior steps whose output it requires in a \texttt{needs} array:

\begin{verbatim}
{"steps": [
  {"description": "...", "toolName": "search_web",
   "toolArgs": {...}, "needs": []},
  {"description": "...", "toolName": "search_web",
   "toolArgs": {...}, "needs": []},
  {"description": "...", "toolName": "analyze_data",
   "toolArgs": {...}, "needs": [0, 1]},
  {"description": "...", "toolName": "write_report",
   "toolArgs": {...}, "needs": [2]}
]}
\end{verbatim}

\noindent This models the plan as a DAG.
The dependency metadata is produced by the planner at zero additional LLM cost---it is part of the same unified planning call (\S\ref{sec:prompt-opt}).
At synthesis time, Joule computes the \emph{sink-node reachability closure}: identify all terminal steps (no dependents), walk backwards via BFS, and include only the union of these closures.

\subsubsection{The Binary Pruning Problem}

A na\"{i}ve implementation prunes all steps outside the reachability closure.
Our ablation (\S\ref{sec:pruning-ab}) shows this causes 0.40--0.57 quality-point regressions on complex tasks.
The cause: LLM planners do not reliably declare \emph{semantic} dependencies.
A step that ``evaluates source credibility'' conceptually influences a later ``write recommendation'' step without producing tokens that appear in the latter's arguments.
Output-signature taint tracking~(\S\ref{sec:taint}) catches data-flow dependencies but not semantic influence.

\subsubsection{Confidence-Aware Pruning}

Instead of binary pruning, Joule computes a \emph{prune confidence score} for each step outside the reachability closure.
A step is pruned only when its score meets or exceeds a dynamically set threshold $\tau_{\text{prune}}$ (see \S\ref{sec:adaptive}).

The score combines three signals, each normalized to [0, 1] where 1.0 = confident it is safe to prune:

\begin{enumerate}[leftmargin=2em,itemsep=2pt]
  \item \textbf{Taint safety} (weight~1.0): If taint tracking detected and repaired a data-flow edge \emph{from} this step to a live step, the planner under-declared at least one dependency here.
        Score = 0.5 (cautious); otherwise 1.0 (no evidence of undeclared data-flow).

  \item \textbf{Lexical relevance} (weight~1.5$\times$ sensitivity): Joule computes the overlap rate between key terms in the step's output and the original task description.
        High overlap suggests the step's output is thematically relevant to the synthesis goal.
        Score = $\max(0,\; 1 - 1.5 \cdot r_{\text{overlap}})$ where $r_{\text{overlap}}$ is the fraction of output terms appearing in the task description.

  \item \textbf{Output richness}: Short outputs (e.g., single-token confirmations) carry little information and are safer to prune than rich multi-sentence outputs.
        Score = 0.90 if $|output| > 200$ chars (richness penalty), else 1.0.
\end{enumerate}

\noindent Combined: $\text{pruneConf}(s) = \text{taintSafety}(s) \cdot \text{lexicalSafety}(s) \cdot \text{richnessFactor}(s)$.
Step $s$ is pruned iff $\text{pruneConf}(s) \geq \tau_{\text{prune}}$.

\subsubsection{Output-Signature Taint Tracking}
\label{sec:taint}

After all steps complete, the runtime extracts \emph{distinctive tokens} from each step's actual output---URLs, file paths, UUIDs, hex hashes, email addresses, long quoted strings---and checks whether any downstream step's \texttt{toolArgs} contains these tokens.
A match constitutes an undeclared data-flow dependency; the edge is automatically inserted before context selection is applied.
Signature classes are designed for zero false positives (Appendix~\ref{app:taint}).

\subsection{Adaptive Controller}
\label{sec:adaptive}

The adaptive controller dynamically sets $\tau_{\text{prune}}$ based on four runtime signals observable at zero additional LLM cost:

\begin{enumerate}[leftmargin=2em,itemsep=2pt]
  \item \textbf{Repair rate}: fraction of dependency edges that were not declared by the planner and had to be repaired by taint tracking. High repair rate signals systematic planner under-declaration.
  \item \textbf{Plan length}: longer plans have more potential for undeclared semantic dependencies between non-adjacent steps.
  \item \textbf{Token budget pressure}: fraction of the task's token budget already consumed. High budget pressure shifts the trade-off toward energy savings.
  \item \textbf{Task complexity}: the planner's estimated complexity (0--1). High complexity tasks benefit from more conservative context selection.
\end{enumerate}

\noindent The threshold is computed as:
\[
  \tau_{\text{prune}} = \tau_{\text{base}}
    + \underbrace{0.15 \cdot \mathbf{1}[\text{repair\_rate} > 0.25]}_{\text{uncertainty penalty}}
    + \underbrace{0.08 \cdot \mathbf{1}[\text{plan\_len} \geq 4]}_{\text{length penalty}}
    + \underbrace{0.05 \cdot \mathbf{1}[\text{complexity} \geq 0.6]}_{\text{complexity penalty}}
    - \underbrace{0.25 \cdot \mathbf{1}[\text{budget\_fraction} \geq 0.80]}_{\text{budget relief}}
\]
where $\tau_{\text{base}} = 0.75$, with the result clamped to $[0.40, 0.95]$.
Under normal conditions, the threshold is 0.75: Joule prunes a step only when 75\%+ confident it is irrelevant.
Under high repair rate and long plan (severe uncertainty), the threshold rises to 0.98: almost nothing is pruned.
Under budget pressure, the threshold drops to 0.50: prune more aggressively to stay within the token ceiling.

The adaptive controller also issues an \emph{escalation recommendation}: escalate to LLM tier when complexity $\geq 0.7$ \emph{and} prior step failures have occurred (the SLM tier was insufficient).
This is a recommendation to the routing layer, not a hard override.

This architecture is \emph{not} a learned model---it is a principled rule set with configurable weights.
Online learning via a bandit-style controller that observes quality outcomes and adjusts weights is left for future work.

\subsection{Adaptive Prompt Optimization}
\label{sec:prompt-opt}

Three optimizations address planning overhead.

\paragraph{Unified Planning.}
Replaces four sequential LLM calls with a single structured JSON response:

\begin{verbatim}
{"goal": "...", "complexity": 0.0-1.0,
 "confidence": 0.0-1.0,
 "steps": [{"description": "...",
             "toolName": "...",
             "toolArgs": {...},
             "needs": [<prior step indices>]}],
 "directAnswer": "..."}
\end{verbatim}

\noindent The \texttt{needs} array is the key schema extension that enables dependency context selection at zero additional LLM cost.
Unified planning collapses 4 round-trips to 1, eliminating 3 network round-trips and their associated system prompt tokens per task.

\paragraph{Slim Prompt Routing.}
The full planning system prompt ($\sim$2,900 tokens, including all tool descriptions) is replaced by a $\sim$50-token slim prompt for tasks identified as unlikely to require tools (no action-intent keywords, description $\leq$300 chars).
This is a 58$\times$ reduction in system prompt tokens for simple tasks.

\paragraph{Direct Answer Shortcut.}
When the unified plan returns a non-empty \texttt{directAnswer} (complexity $\leq$0.5, empty step list), Joule returns it directly without execution or synthesis.
This eliminates 1+ additional LLM round-trips per low-complexity task.

\subsection{SLM-First Routing}
\label{sec:routing}

Every LLM call passes through a model router applying a five-rule priority cascade:

\begin{enumerate}[leftmargin=2em,itemsep=2pt]
  \item \textbf{Purpose guard}: Classification and verification always use SLM.
  \item \textbf{Escalation budget exhausted}: Hard constraint to SLM.
  \item \textbf{Energy ceiling}: When energy proxy remaining $<$ 0.01\,Wh, force SLM.
  \item \textbf{Complexity threshold}: If complexity $> \tau_c = 0.7$, escalate to LLM.
  \item \textbf{Confidence threshold}: If prior step confidence $< \tau_s = 0.6$, escalate.
\end{enumerate}

\noindent Default: SLM. Routing to SLM vs.\ LLM yields a 4--5$\times$ energy proxy reduction per token.
The adaptive controller's escalation recommendation (\S\ref{sec:adaptive}) feeds into the routing context for Rule~4 as an additional complexity signal.

\subsection{7-Dimensional Budget Enforcement}
\label{sec:budget}

Each task is assigned a budget envelope across seven dimensions (Table~\ref{tab:budget-presets}).
After every LLM call, tool invocation, and escalation, the manager checks:
$\forall d: \text{consumed}(d) \leq \text{ceiling}(d)$.
Violation halts execution immediately (hard synchronous check, not a warning).

\begin{table}[h]
\centering
\caption{Budget envelope dimensions and preset values.}
\label{tab:budget-presets}
\begin{tabular}{llrrr}
\toprule
Dimension & Unit & Low & Medium & High \\
\midrule
Tokens         & count   & 4{,}000   & 16{,}000  & 100{,}000 \\
Cost           & USD     & \$0.01    & \$0.10    & \$1.00    \\
Latency        & ms      & 10{,}000  & 30{,}000  & 300{,}000 \\
Tool calls     & count   & 3         & 10        & 40        \\
Escalations    & count   & 0         & 1         & 5         \\
Energy proxy   & Wh      & 0.005     & 0.05      & 0.5       \\
Carbon         & gCO$_2$ & 0.002     & 0.02      & 0.2       \\
\bottomrule
\end{tabular}
\end{table}

\noindent Budget enforcement also \emph{steers} routing: exhausted escalation allowance forces SLM (Rule~2); near-ceiling energy forces SLM (Rule~3).
This bidirectional coupling distinguishes Joule from a simple hard stop.

% ─────────────────────────────────────────────────────────────────────
\section{Experimental Setup}
\label{sec:setup}

\paragraph{Benchmark.}
50 tasks spanning five categories (10 each):
\emph{QA} (factual, no tool use), \emph{summarization} (condense text),
\emph{research} (multi-source synthesis, 2--3 tool calls),
\emph{code generation} (write/explain code),
\emph{multi-step} (chained tool pipelines, 2--3 steps).
3 runs per task; all reported values are mean $\pm$ stddev over 150 executions.

\paragraph{Baselines.}
(1)~\textbf{CrewAI}~\cite{crewai}: single-agent, role-based prompting.
(2)~\textbf{LangChain}~\cite{langchain}: ReAct-style tool-use agent.
(3)~\textbf{Direct}: a single \texttt{gpt-4o-mini} call with no planning, included as a minimal-overhead reference for simple tasks.
All configurations use \texttt{gpt-4o-mini} to ensure energy proxy comparisons reflect token counts, not model selection.

\paragraph{Energy proxy.}
$E = (n_{\text{in}} \cdot 0.30 + n_{\text{out}} \cdot 1.20) / 10^6$\,Wh for \texttt{gpt-4o-mini}.
We explicitly call this an energy \emph{proxy}: it models inference cost proportional to token consumption and per-model efficiency ratios from~\cite{luccioni2023power}.
Actual hardware power consumption varies by deployment; we do not claim these are direct Wh measurements.

\paragraph{Quality.}
LLM-as-judge scoring on a 1--5 scale, standardized rubric, applied uniformly to all frameworks.

\paragraph{Ablation configurations.}
\textbf{full} (all mechanisms); \textbf{no-prompt-opt} (separate spec/classify/plan/critique calls); \textbf{no-routing} (all calls forced to LLM); \textbf{no-pruning} (all step results passed to synthesis, equivalent to binary pruning with $\tau_{\text{prune}} = 0$).

% ─────────────────────────────────────────────────────────────────────
\section{Evaluation}
\label{sec:eval}

\subsection{Overall Comparison}
\label{sec:overall}

% Auto-generated LaTeX tables for Joule paper
% Generated: 2026-03-12 (adaptive controller benchmark — confidence-aware pruning)
% Frameworks: joule-full, joule-no-pruning, joule-no-prompt-opt, joule-no-routing, crewai, langchain
% Runs per task: 3 (150 total executions)

% ── Table 1: Overall comparison ──────────────────────────────────────

\begin{table}[h]
\centering
\caption{Overall comparison: Joule vs.\ CrewAI vs.\ LangChain across 50 tasks, 3 runs per task (150 executions). Values are mean $\pm$ stddev. Joule is 2.4$\times$ faster than both baselines with equivalent success rate; per-category energy analysis (Table~\ref{tab:category-energy}) shows Joule reduces inference cost by 28\% on research and 30\% on multi-step tasks.}
\label{tab:overall-comparison}
\begin{tabular}{l r r r}
\toprule
Metric & Joule & CrewAI & LangChain \\
\midrule
Success Rate         & \textbf{100.0\%} & 100.0\%  & 99.3\%   \\
Avg Prompt Tokens    & 2{,}288 $\pm$ 844   & 441 $\pm$ 491    & 619 $\pm$ 719    \\
Avg Completion Tokens& 256 $\pm$ 162    & 683 $\pm$ 588    & 678 $\pm$ 596    \\
Avg LLM Calls/Task   & \textbf{1.5 $\pm$ 0.7} & 1.4 $\pm$ 0.5 & 1.8 $\pm$ 0.7 \\
Avg Energy (Wh)      & 0.000994 $\pm$ 0.000432 & \textbf{0.000951 $\pm$ 0.000849} & 0.000998 $\pm$ 0.000925 \\
Avg Latency (ms)     & \textbf{5{,}560 $\pm$ 3{,}596} & 13{,}143 $\pm$ 11{,}054 & 12{,}797 $\pm$ 11{,}554 \\
Avg Quality (1--5)   & 4.61 $\pm$ 0.71  & 4.86 $\pm$ 0.29  & 4.89 $\pm$ 0.28  \\
Budget Violations    & \textbf{0 / 150} & N/A & N/A \\
\bottomrule
\end{tabular}
\end{table}

% ── Table 2: Per-category energy ─────────────────────────────────────

\begin{table}[h]
\centering
\caption{Per-category energy consumption (Wh), mean $\pm$ stddev, 3 runs per task.
Joule reduces energy by 28\% and 30\% on research and multi-step tasks respectively vs.\ CrewAI.
For simple tasks (QA, summarization, code generation), baselines' lack of planning overhead gives them an energy advantage.
\textbf{Bold} marks the lowest energy per row.}
\label{tab:category-energy}
\begin{tabular}{l r r r}
\toprule
Category & Joule & CrewAI & LangChain \\
\midrule
QA           & 0.000489 $\pm$ 0.000026 & \textbf{0.000063 $\pm$ 0.000090} & 0.000060 $\pm$ 0.000081 \\
Summarization& 0.000690 $\pm$ 0.000218 & \textbf{0.000254 $\pm$ 0.000140} & 0.000252 $\pm$ 0.000143 \\
Research     & \textbf{0.001289 $\pm$ 0.000138} & 0.001800 $\pm$ 0.000064 & 0.001955 $\pm$ 0.000227 \\
Code Gen.    & 0.001029 $\pm$ 0.000432 & \textbf{0.000548 $\pm$ 0.000142} & 0.000550 $\pm$ 0.000166 \\
Multi-step   & \textbf{0.001472 $\pm$ 0.000126} & 0.002092 $\pm$ 0.000154 & 0.002202 $\pm$ 0.000319 \\
\bottomrule
\end{tabular}
\end{table}

% ── Table 3: Per-category latency ────────────────────────────────────

\begin{table}[h]
\centering
\caption{Per-category mean latency (ms). Joule is 3.9$\times$ faster on research tasks and 2.8$\times$ faster on multi-step tasks vs.\ both baselines. On simple tasks (QA, summarization), Joule's planning step adds moderate overhead. \textbf{Bold} marks the lowest latency per row.}
\label{tab:category-latency}
\begin{tabular}{l r r r}
\toprule
Category & Joule & CrewAI & LangChain \\
\midrule
QA            & 1{,}797 & \textbf{1{,}608} & 1{,}045 \\
Summarization & 3{,}708 & \textbf{3{,}597} & -- \\
Research      & \textbf{6{,}478} & 25{,}027 & 27{,}834 \\
Code Gen.     & \textbf{6{,}372} & 8{,}948 & 8{,}108 \\
Multi-step    & \textbf{9{,}445} & 26{,}534 & 28{,}863 \\
\bottomrule
\end{tabular}
\end{table}

% ── Table 4: Ablation study ──────────────────────────────────────────

\begin{table}[h]
\centering
\caption{Ablation study: contribution of each Joule mechanism (50 tasks, 3 runs each).
Disabling prompt optimization (\emph{no-prompt-opt}) multiplies energy by 4.0$\times$ and LLM calls by 3.5$\times$, making it the dominant mechanism.
Adaptive pruning (\emph{full} vs.\ \emph{no-pruning}) raises quality by 0.27 points at 6\% energy overhead vs.\ no-pruning,
confirming it is a quality-preservation mechanism rather than an energy-savings mechanism.
\textbf{Bold} marks the best value in each column among Joule configurations.}
\label{tab:ablation}
\begin{tabular}{l r r r r r}
\toprule
Config & Energy (Wh) & Tokens & LLM Calls & Quality & Success \\
\midrule
\textbf{full}    & 0.000994 $\pm$ 0.000432 & 2{,}544 $\pm$ 965 & \textbf{1.5} & 4.61 $\pm$ 0.71 & \textbf{100.0\%} \\
no-pruning       & \textbf{0.000936 $\pm$ 0.000447} & \textbf{2{,}341 $\pm$ 974} & \textbf{1.5} & \textbf{4.69 $\pm$ 0.58} & 99.3\% \\
no-routing       & 0.001013 $\pm$ 0.000454 & 2{,}561 $\pm$ 981 & \textbf{1.5} & \textbf{4.69 $\pm$ 0.53} & \textbf{100.0\%} \\
no-prompt-opt    & 0.003929 $\pm$ 0.000530 & 11{,}542 $\pm$ 1{,}489 & 5.3 & 4.53 $\pm$ 0.67 & 97.3\% \\
\midrule
\textit{CrewAI}  & \textit{0.000951} & \textit{1{,}124} & \textit{1.4} & \textit{4.86} & \textit{100.0\%} \\
\textit{LangChain} & \textit{0.000998} & \textit{1{,}292} & \textit{1.8} & \textit{4.89} & \textit{99.3\%} \\
\bottomrule
\end{tabular}
\end{table}

% ── Table 5: Pruning A/B ─────────────────────────────────────────────

\begin{table}[h]
\centering
\caption{Dependency pruning A/B: confidence-aware adaptive pruning (\emph{full}) vs.\ no-pruning baseline.
The adaptive controller sets pruning rate to 0\% on all categories, correctly identifying all step outputs as high-relevance.
This eliminates the multi-step quality regression entirely ($\Delta$Q\,=\,0.00 on multi-step) and reduces the research gap from 0.57 to 0.36 points vs.\ binary pruning (prior work).
Negative E-saved values ($<$2\%) reflect the negligible overhead of retaining additional context.}
\label{tab:pruning-ab}
\begin{tabular}{l r r r r r r r}
\toprule
Category & \multicolumn{2}{c}{Quality (1--5)} & $\Delta$Q & \multicolumn{2}{c}{Energy (Wh)} & E-saved & Prune\% \\
\cmidrule(lr){2-3}\cmidrule(lr){5-6}
 & Adaptive & No-Prune & & Adaptive & No-Prune & (\%) & \\
\midrule
QA            & 5.00 & 5.00 & \phantom{$-$}0.00  & 0.000489 & 0.000431 & $-1.4$ & 0.0\% \\
Summarization & 5.00 & 5.00 & \phantom{$-$}0.00  & 0.000690 & 0.000623 & $-1.1$ & 0.0\% \\
Research      & 3.67 & 4.03 & $-$0.36 & 0.001289 & 0.001276 & $-1.0$ & 0.0\% \\
Code Gen.     & 5.00 & 5.00 & \phantom{$-$}0.00  & 0.001029 & 0.000903 & $-1.4$ & 0.0\% \\
Multi-step    & 4.40 & 4.40 & \phantom{$-$}0.00  & 0.001472 & 0.001459 & $-0.9$ & 0.0\% \\
\bottomrule
\end{tabular}
\end{table}


Table~\ref{tab:overall-comparison} reports head-to-head results.
Joule achieves comparable overall energy proxy (0.000994\,Wh vs.\ 0.000951 for CrewAI and 0.000998 for LangChain), while being \textbf{2.4$\times$ faster} (5,560\,ms vs.\ 13,143\,ms and 12,797\,ms).
The per-category breakdown (Table~\ref{tab:category-energy}) reveals where the efficiency story lives: Joule's planning overhead adds cost on simple tasks, but saves 28--30\% energy on complex tasks.

The latency advantage is structural: both baselines use iterative observe-think-act loops whose latency scales with output generation; Joule's unified planning call returns a structured plan quickly, with execution stages focused and parallelizable.

Quality scores are 4.61 for Joule vs.\ 4.86 (CrewAI) and 4.89 (LangChain).
The gap on QA and code generation is small; the remaining gap is concentrated in \emph{research} (\S\ref{sec:pruning-ab}).
Budget enforcement holds: zero violations across all 150 executions.

\subsection{Per-Category Analysis}
\label{sec:percategory}

Table~\ref{tab:category-energy} shows per-category energy proxy and Table~\ref{tab:category-latency} shows latency.

\paragraph{Research tasks (Joule wins on complex tasks).}
Joule: 0.001289\,Wh.
CrewAI: 0.001800 (+40\%). LangChain: 0.001955 (+52\%).
Savings: \textbf{28\%} vs.\ CrewAI, \textbf{34\%} vs.\ LangChain.
These arise from unified planning eliminating multi-call overhead, and focused per-step context keeping individual prompt sizes bounded.
Latency: Joule 6,478\,ms vs.\ 25,027\,ms (CrewAI) and 27,834\,ms (LangChain)---a \textbf{3.9$\times$ speedup}.

\paragraph{Multi-step tasks (Joule wins on complex tasks).}
Joule: 0.001472\,Wh.
CrewAI: 0.002092 (+42\%). LangChain: 0.002202 (+50\%).
Savings: \textbf{30\%} vs.\ CrewAI, \textbf{33\%} vs.\ LangChain.
Latency: Joule 9,445\,ms vs.\ 26,534 and 28,863\,ms---\textbf{2.8$\times$ speedup}.

\paragraph{QA and summarization (planning overhead dominates).}
For simple QA: Joule 0.000489\,Wh vs.\ CrewAI 0.000063 and LangChain 0.000060.
Joule consumes $\sim$8$\times$ more energy on simple factual questions.
This is expected and unavoidable in a plan-then-execute architecture: even with slim prompts and direct-answer shortcuts (which fire on 100\% of QA tasks), the planning call itself incurs tokens that a direct LLM call avoids.
Teams deploying Joule on QA-dominated workloads should use the direct-call pattern.
For mixed workloads where complex tasks are the common case, the research/multi-step savings dominate in aggregate.

\paragraph{Code generation.}
Joule: 0.001029\,Wh vs.\ 0.000548 (both baselines).
Planning overhead is not justified by the single-completion nature of most code generation tasks in our benchmark.
Code generation tasks requiring iterative tool use (compile-debug loops) would be expected to show Joule advantages not captured here.

\subsection{Ablation Study}
\label{sec:ablation}

Table~\ref{tab:ablation} reports ablation results.

\paragraph{Prompt optimization is the dominant mechanism.}
\textbf{no-prompt-opt}: 0.003929\,Wh (4.2$\times$ more), 11,023 prompt tokens (5.1$\times$), 5.3 LLM calls (3.5$\times$), 13,799\,ms latency (2.9$\times$ slower), 97.3\% success (vs.\ 100\%).
The 4.2$\times$ energy multiplier is the paper's strongest finding: it quantifies how wasteful the na\"{i}ve multi-call planning pattern is, independently of model choice.
Existing frameworks that use this pattern are paying this overhead on every task.

\paragraph{Routing provides secondary savings.}
\textbf{no-routing}: 0.001013\,Wh (+9\%).
The modest impact reflects that our benchmark tasks stayed within SLM confidence thresholds; no LLM escalations were triggered in the full configuration.
The routing benefit would be more pronounced on workloads with tasks genuinely requiring LLM-tier reasoning.

\paragraph{Adaptive pruning: quality preservation at modest cost.}
\textbf{no-pruning}: 0.000936\,Wh vs.\ full 0.000994\,Wh (6\% overhead for retaining additional context).
The confidence-aware adaptive controller sets prune rate to 0\% on all task categories in this benchmark, correctly determining that step outputs are high-relevance and should not be pruned.
The quality improvement over the previous binary-pruning approach (+0.12 overall, +0.20 on research, +0.40 on multi-step) demonstrates that the adaptive threshold is functioning as designed.
Per-category effects are analyzed in \S\ref{sec:pruning-ab}.

\subsection{Context Selection A/B}
\label{sec:pruning-ab}

Table~\ref{tab:pruning-ab} compares \textbf{full} (confidence-aware context selection with adaptive controller, $\tau_{\text{base}}=0.75$) vs.\ \textbf{no-pruning} (all step results unconditionally included).

\paragraph{Adaptive controller behavior.}
The adaptive controller set prune rate to \textbf{0\%} on every task category.
This is the expected outcome: research and multi-step step outputs exhibit high lexical overlap with the task description (overlapRate $\approx 0.35$--0.50), driving lexical safety scores of $\approx 0.40$--0.50, well below the $\tau_{\text{prune}} = 0.75$ threshold.
The controller correctly identifies these as high-relevance steps and retains them.

\paragraph{Quality results.}
Compared to the prior binary-pruning implementation:
\begin{itemize}[leftmargin=2em,itemsep=2pt]
  \item Multi-step: quality 4.40 (vs.\ 4.00 with binary pruning, 4.40 with no-pruning) --- \textbf{full quality recovery}.
  \item Research: quality 3.67 (vs.\ 3.47 with binary pruning, 4.03 with no-pruning) --- \textbf{partial recovery}, gap reduced from 0.57 to 0.36.
  \item QA, summarization, code generation: 5.00 / 5.00 / 5.00 --- no change.
\end{itemize}
Overall quality improves from 4.49 (binary pruning) to 4.61 (adaptive controller).

\paragraph{Energy results.}
Energy overhead vs.\ no-pruning is negligible on complex tasks: $-$1.0\% on research, $-$0.9\% on multi-step (negative means full uses marginally more energy than no-pruning due to retaining additional context in synthesis).
Values are within run-to-run variance; no meaningful energy difference exists at 0\% prune rate.

\paragraph{Discussion.}
The 0\% prune rate confirms the adaptive controller's conservatism mechanism is functioning: it identifies high-relevance outputs and refuses to prune them.
This is the correct behavior---the quality improvements over binary pruning validate that the prior 0.40--0.57 regressions were caused by incorrect pruning decisions.
The residual 0.36-point research gap vs.\ no-pruning reflects that research tasks have genuine semantic dependencies that neither taint tracking nor lexical overlap fully captures.

The path to closing this remaining gap is \emph{improved planner dependency declaration}: if the planner explicitly declares all semantic dependencies in the \texttt{needs} array, taint tracking and lexical scoring become precise, and high-confidence pruning of genuinely irrelevant steps becomes achievable.
Chain-of-thought dependency reasoning in the planner or fine-tuning on dependency annotation are the most direct paths to this improvement.

\subsection{Discussion}
\label{sec:discussion}

\paragraph{Where Joule provides value.}
Research and multi-step tasks (28--30\% energy savings, 2.8--3.9$\times$ speedup).
Any workload with multi-call planning pipelines (4.0$\times$ energy multiplier from prompt optimization alone).
Deployments requiring hard resource guarantees (zero violations across all runs).

\paragraph{Where Joule adds overhead.}
Simple QA and single-completion tasks.
The planning infrastructure is not free for trivial requests.
The break-even point is approximately any task requiring 2+ tool calls.

\paragraph{Measurement validity.}
Three runs per task provide variance estimates.
All frameworks use identical \texttt{gpt-4o-mini} models---energy differences are driven by token counts, not model selection.
Joule shows lower per-task energy variance (stddev 0.000432) than baselines (0.000849 for CrewAI, 0.000925 for LangChain), reflecting the predictability of a structured plan-then-execute pipeline.
Energy values are token-weighted proxies; direct hardware power measurement is a direction for future work.

% ─────────────────────────────────────────────────────────────────────
\section{Related Work}
\label{sec:related}

\paragraph{Energy-efficient inference.}
Model-level efficiency through quantization~\cite{dettmers2023qlora}, speculative decoding~\cite{leviathan2023fast}, and serving optimization~\cite{kwon2023efficient} reduces per-token cost.
Luccioni et al.~\cite{luccioni2023power} provide empirical per-model measurements.
These approaches are complementary: Joule operates at the agent runtime layer above the model, reducing token counts rather than per-token costs.

\paragraph{Agent frameworks.}
CrewAI~\cite{crewai}, LangChain~\cite{langchain}, AutoGen~\cite{autogen} focus on capability, not efficiency.
LangChain's ReAct~\cite{yao2022react} uses multi-call observe-think-act loops; Joule's unified planning achieves equivalent decomposition in one call.
No prior framework models plan dependency structure for context selection.

\paragraph{Context compression.}
RAG~\cite{lewis2020rag} and prompt compression~\cite{jiang2023llmlingua} reduce context volume by filtering or compressing retrieved content.
Joule's context selection is structurally different: it uses the execution plan's dependency graph---not content similarity or compression---to determine relevance.
No additional LLM calls are required; the pruning cost is a CPU-side BFS over a graph with at most tens of nodes.

\paragraph{Budget and resource management.}
Budget-aware scheduling exists in cloud computing~\cite{xiao2013dynamic} and hyperparameter search~\cite{liaw2018tune}.
Token budget management in long-context inference~\cite{shi2024context} targets single-call efficiency.
Joule's 7-dimensional budget model spanning multi-step task execution is, to our knowledge, the first such system for agent runtimes.

% ─────────────────────────────────────────────────────────────────────
\section{Conclusion}
\label{sec:conclusion}

We presented Joule, an execution-aware runtime for multi-step LLM agent workflows.
Its central contribution is confidence-aware dependency context selection: modeling execution plans as DAGs, scoring candidate-prune steps via taint tracking, lexical relevance, and output richness, and using an adaptive controller to set the pruning conservatism threshold based on observed planning uncertainty and budget pressure.
Empirically, the adaptive controller correctly identifies high-relevance step outputs and retains them, fully recovering the quality regression caused by binary pruning (multi-step quality 4.00 $\to$ 4.40) while reducing the research gap from 0.57 to 0.36 points.

The dominant empirical finding is that \textbf{prompt optimization---not context selection---is the largest avoidable overhead}: unifying four planning calls into one reduces per-task energy proxy by 4.0$\times$ and LLM call count by 71\%.
This implies that existing frameworks using multi-call planning pipelines are paying a 4$\times$ overhead on every task.

On the workloads where agent frameworks provide value (research synthesis, multi-step automation), Joule achieves 28--30\% lower energy proxy and 2.8--3.9$\times$ lower latency than CrewAI and LangChain.
Budget enforcement provides formal per-task resource guarantees unavailable in prior frameworks.

Future work includes: (1)~improving planner dependency declaration quality to increase pruning precision; (2)~online learning for the adaptive controller's threshold weights using execution quality feedback; (3)~direct hardware energy measurement to validate the proxy model.

\begin{acks}
The author thanks the University of Illinois Chicago systems group for compute access.
\end{acks}

\bibliographystyle{ACM-Reference-Format}
\bibliography{joule-refs}

% ─────────────────────────────────────────────────────────────────────
\appendix

\section{Energy Proxy Derivation}
\label{app:energy}

Coefficients in Table~\ref{tab:energy-tiers} are from Luccioni et al.~\cite{luccioni2023power}.
Input prefill and output decoding are metered separately because decoding is memory-bandwidth bound and 3--5$\times$ more energy-intensive per token than prefill.
For \texttt{gpt-4o-mini}: $\alpha_{\text{in}} = 0.30$\,Wh/M, $\alpha_{\text{out}} = 1.20$\,Wh/M.
These are proxies; actual values depend on datacenter efficiency, GPU utilization, and batching.

For local models (\texttt{llama3.2:3b} via Ollama): 65\,W TDP CPU at 20\,tok/s $\Rightarrow$ 3.25 Wh/M tokens (all tokens, undifferentiated by direction).
Direct measurement is future work.

\section{Taint Tracking Signature Classes}
\label{app:taint}

\begin{table}[h]
\centering
\caption{Output signature classes for taint tracking validation. All classes designed for zero false positives. Cap: 50 signatures per step.}
\label{tab:sig-classes}
\begin{tabular}{llc}
\toprule
Class & Example & Min Len \\
\midrule
URL & \texttt{https://api.example.com/v2/} & 10 \\
File path & \texttt{/home/user/data/report.csv} & 8 \\
UUID & \texttt{550e8400-e29b-41d4-a716-...} & 36 \\
Hex hash & \texttt{a1b2c3d4e5f6a7b8c9d0...} & 16 \\
Email & \texttt{user@example.com} & 6 \\
Distinctive string & \texttt{``Q3 Revenue Analysis 2025''} & 20 \\
\bottomrule
\end{tabular}
\end{table}

\section{Adaptive Controller Threshold Formula}
\label{app:adaptive}

The threshold $\tau_{\text{prune}}$ is clamped to $[0.40, 0.95]$ after adjustment:

\begin{align*}
\tau_{\text{prune}} &= 0.75 \\
  &+ 0.15 \cdot \mathbf{1}[\text{repair\_rate} > 0.25] \\
  &+ 0.08 \cdot \mathbf{1}[\text{plan\_len} \geq 4] \\
  &+ 0.05 \cdot \mathbf{1}[\text{complexity} \geq 0.6] \\
  &- 0.25 \cdot \mathbf{1}[\text{budget\_fraction} \geq 0.80]
\end{align*}

Example values: no pressure signals $\to$ 0.75; high repair rate + long plan $\to$ 0.98; budget pressure only $\to$ 0.50; all signals $\to$ 0.73.

\end{document}
